\chapter{Memory}\label{ch:memory}

Loads and stores are per lane. Each instruction names a
base (a device binding, threadgroup memory, the imageblock, or a register-held pointer), an index (a register or an
immediate) and a byte displacement (\cref{fig:address-spaces}). The opcode fixes the element width; there is no width field.
Tensor loads are per-lane gathers through the same load and store machinery and the same eight-slot scoreboard (MM~\mm{7};
\cref{sec:tensor-memory-path}).

The hardware does not bounds-check memory accesses, and it does not check that a loaded value has arrived
(\cref{sec:scoreboard}). With a wrong binding rank (\cref{sec:address-computation}), or a consumer that does not wait
on its load, the dispatch completes with status 0 and returns wrong data.

\leadhead{Notation.}

\begin{itemize}
\item \emph{Rn}: 32-bit general register n, named here by role ("R32, the per-lane index register").
\item \emph{\opn{0}} and \emph{\opn{4}}: the decoder's tags for the uniform register file and for the message staging file that texture
  messages read.
\item \emph{Binding expression}: how the decoder prints a base, \code{bin(op0, const(K), 8)} for a device base and
  \code{bin(op0, const(K), 2)} for a threadgroup base. K is the binding constant and the last number is the base's scale.
\item \emph{Binding rank}: a buffer's position in the binding list the program image declares (\cref{sec:address-computation}).
\item \emph{Operand 1}: the control word after a memory instruction's first operand. It carries the slot a producer fills
  and the instruction's own wait mask (\cref{sec:scoreboard}).
\item \emph{Element code}: the access-width field of tensor and vector loads: 1, 2, 4, 8 or 16 bytes (MM~\mm{7}, MM~\mm{25.37}).
\item \emph{Immediate-address form}: a load or store whose index and displacement are both immediates (\cref{sec:immediate-address-vector-forms}).
\item "Below Metal" means a pure IOGPU submission with no AGXMetal driver loaded
  (\code{docs/g17-first-dispatch-trials.md}; interface in \cref{ch:submission-below-metal}).
\end{itemize}

% Where a load or store can point, and the limit that governs each space.
\begin{figure}[htbp]
\centering
\begin{tikzpicture}[node distance=3mm,
    space/.style={box, draw=ink!55, fill=ink!5, text width=50mm, align=left, font=\scriptsize,
      minimum height=11mm},
    lim/.style={tag, font=\scriptsize, align=left, text width=40mm}]

  \node[space] (dev)
    {\textbf{Device memory}\\[-1pt]named by binding rank; the decoder prints
     \code{bin(op0, const(K), 8)}};
  \node[lim, right=5mm of dev] (devl)
    {No bounds check. A wrong rank returns wrong data with status 0.};

  \node[space, below=of dev] (tg)
    {\textbf{Threadgroup memory}\\[-1pt]\code{bin(op0, const(K), 2)}, shared by the threadgroup};
  \node[lim, right=5mm of tg] (tgl)
    {Apple's compiler refuses past 32{,}768\,B. 57{,}344\,B per core, inferred from a step in
     run time.};

  \node[space, below=of tg] (ib)
    {\textbf{Imageblock}\\[-1pt]per-threadgroup tile, addressed by an (x, y) register};
  \node[lim, right=5mm of ib] (ibl)
    {x in the low half, y in the high half. The emulator refuses coordinates outside the tile;
     hardware completed one that was.};

  \node[space, below=of ib] (pool)
    {\textbf{Constant pool}\\[-1pt]compile-time data in metadata slot 13};
  \node[lim, right=5mm of pool] (pooll)
    {Read as \code{bin(op0, const(K), W)} --- the same syntax as a uniform.};

  \node[space, below=of pool] (uni)
    {\textbf{Uniform file}\\[-1pt]written per dispatch by the constant program};
  \node[lim, right=5mm of uni] (unil)
    {Indistinguishable from a pool entry in the operand alone; confusing the two gave six silent
     wrong answers.};

  \node[space, below=of uni] (tex)
    {\textbf{Textures}\\[-1pt]reached by message, operands staged in the \opn{4}};
  \node[lim, right=5mm of tex] (texl)
    {The staging write names a scoreboard slot; the message waits on it.};

  \draw[decorate, decoration={brace, amplitude=5pt}, line width=0.6pt, color=ink!70]
    ([xshift=-2mm]dev.north west) -- ([xshift=-2mm]tex.south west)
    node[midway, left=6pt, tag, text width=26mm, align=right]
      {one load or store names exactly one of these, plus an index and a displacement};
\end{tikzpicture}
\caption{Where an access can point. The opcode fixes the element width; there is no width field, and the
hardware checks neither the bound nor whether a loaded value has arrived.}
\label{fig:address-spaces}
\end{figure}


\section{Address computation}\label{sec:address-computation}

A register-indexed access forms its address as

\begin{lstlisting}
address = base + index_register * element_width + displacement        (width and displacement in bytes)
\end{lstlisting}

Only the index is scaled. A constant address has no index register, width 1 and a byte
displacement. The immediate-address family (\cref{sec:immediate-address-vector-forms}) uses a different formula.

Two project files give conflicting answers on whether the displacement exists and what it counts.
\code{isa/g17-addressing-rules.json} compiled \code{C[tid] = A[tid + k]} with \code{uint tid} for six element types and found that no operand of the load moved with k: Apple's compiler instead added k to the index with \op{10279}, and k = 1, 2, 4, 8 encoded as immediate 1, 2, 4, 8 for char through long. Its conclusion was "the
load carries NO displacement field" and \code{address = base + (index + displacement) * sizeof}. \code{isa/g17-address-formula.txt}
compiled \code{f[i + 7]}, \code{w[i + 7]} and \code{L[i + 7]} and found a displacement operand carrying 28, 14 and 56: bytes, not
elements. The first conclusion holds only for its unsigned \code{tid + k} sweep (MM~\mm{7}).

On hardware, setting the 128-bit vector load's displacement to 4, 8, 12 and 20 moved the access by exactly that many bytes in 10 of
10 runs (MM~\mm{25.141.18}). The repository's emulator, which models every register-indexed access as
\code{index * element + displacement} in bytes, produces the GPU's output exactly on 316 of 316 and 172 of 172 shipped
dispatches (MM~\mm{25.197}).

A constant address carries its byte offset in the displacement: \code{u[1]} is 4, \code{u[7]} 28, \code{u[64]} 256 for uint
(Apple's compiler emits, \code{isa/g17-address-formula.txt}). For \code{A[tid + k]} with an unsigned 32-bit index, Apple's
compiler puts k into the index with the add immediate (1 to 255) or, at 256 and above, \op{10282} or \op{10283}, and leaves
the displacement 0 (Apple's compiler emits, \code{isa/g17-addressing-rules.json}). The reason is not recorded, so it is not
known whether the displacement can stand in for that add when an unsigned index sum may wrap. The add immediate scatters
its 8-bit value across the word (bits 0--1 at byte1{[}0:2{]}, 2--4 at byte3{[}5:8{]}, 5 at byte5{[}7{]}, 6--7 at
byte8{[}0:2{]}), and this map predicted the encodings of 15 values never seen compiled (decoder, same file).

Apple's compiler emits no index-scale field (\code{isa/g17-address-formula.txt},
\code{isa/g17-addressing-rules.json}): \code{u[i * 4]} emits the same load
as \code{u[i]}, with the multiply in a separate ALU instruction. The opcode carries the access width. char, short and half
use \op{12646} and \op{17193}; int and float use \op{12682} and \op{17229}; long uses \op{12691} and the store \op{17238}
(no glossary entry); int4 uses \op{12709} and \op{17256}. It carries the index width as well: a ulong index selects
\op{12688}, preceded by a 64-bit shift pair that scales the index. A runtime base plus index (\code{A[j + tid]}, j loaded)
is a load of j, a 32-bit add, then the indexed load.

Misalignment does not select a form: an int at byte offset 1 through a char-pointer cast compiles to the same plain 32-bit
load and store; whether the hardware permits a misaligned scalar access is open. A vector load displaced by 4 bytes, so not
16-byte aligned, read all four words from the displaced address (measured, MM~\mm{25.141.18}). No negative displacement was
ever emitted (the compiler folds negative offsets elsewhere), and the field's signedness is unknown.

The address is formed in 64 bits (measured, \code{isa/g17-predication.json}; a load $2^{32}$ bytes past a 256-word
buffer, where a 32-bit wrap would land back inside, returned 0, \cref{sec:out-of-range-accesses-observed}).

A device base's binding constant is K = 4 $\times$ the binding rank, and the rank is the buffer's
position in the list the image declares, not the Metal slot number: a kernel binding only buffer C makes C rank 0
(measured, \code{agxforge/g17/cc.py}; MM~\mm{0.8}). The list order is internal bindings ascending, then user bindings ascending;
reconstructed from the declaration it reproduces Apple's list in 10,964 of 12,044 corpus kernels (91.0\%), never
differing by order alone (Apple's compiler emits, \code{agxforge/g17/resource.py}). A wrong rank makes a store return its fill
value at status 0 (measured, same file).

Below Metal, declaring one unused buffer between A and C changed exactly one code byte, the store's descriptor operand,
from 0x04 to 0x08 (measured, \code{docs/g17-first-dispatch-trials.md}, "Authored device atomic below Metal" and "Input
descriptor address selection" onward). Swapping two 64-bit
entries of the argument block swapped which buffer a load read. When an entry pointed at an unviewed offset of an
allocation, or at a different registered allocation, loads and stores both followed it.

The 128-bit vector load takes its 14-byte form when it carries bit 37 or a
displacement above 255 (decoder, MM~\mm{6}). Tensor loads carry up to 255 bytes in the 12-byte form and 32,767 in the
16-byte form; past 32,767 the offset is folded into a per-tile base register (Apple's compiler emits, MM~\mm{25.114.3}).
This compiler does the same, visible as a kernel growing from 74 to 80 instructions, with no failure out to 1 MiB of
leading dimension (measured, MM~\mm{15}). An identical offset-bearing body was bit-exact at stream positions 1 to 12,
after 0 to 4,096 bytes of padding and at offsets 2 to 40,002 bytes straddling the boundary; within that range there is no
stream-position constraint (233 + 193 dispatches, MM~\mm{15}). Tensor B offsets are measured from 2 to 47,104
bytes (MM~\mm{25.114.3}).

Tensor loads use this formula with an index that ordinary arithmetic computes (measured, MM~\mm{7}). Apple's library computes one per-lane index register,
R32, the element offset of the lane's first element in a row-major tile, and its eight tensor loads read
\code{[R32*2 + disp]} at the byte displacements \{0, 32, 2048, 2080, 4096, 4128, 6144, 6176\}: four 8-row steps of 2,048
bytes times two 16-column steps of 32, each load moving an 8-row by 16-column sub-tile (the full set is in
\code{docs/g17-tensorops-accelerator-recon.md} section 93); the fragment layout follows from this arithmetic (\cref{sec:fragment-layouts}). There is no row-stride
field. A non-default stride changes the mode immediate from 241 to 2289
and moves into the index (decoder, MM~\mm{7}); a stride built as \code{K*2} regardless of element size was wrong, with no
error reported, on 1,024 of 1,024 elements of a 32 $\times$ 32 $\times$ 64 int8 product (measured, MM~\mm{17} rule 6; \cref{sec:tensor-memory-path}).

\section{Load and store forms by width}\label{sec:load-store-forms}

\Cref{tab:memory-forms} lists, by width, the forms this chapter uses; \cref{app:c} lists every memory opcode. Lengths are those
seen in Apple's 6,594-program corpus (\code{isa/g17-wait-mask-census.json}), plus any the cited section measured. The plain
32-bit store's operands are the value (0), operand 1 (1), the index (5) and the index lifetime (6) (\cref{sec:scoreboard}).
The threadgroup and imageblock operand orders are in \cref{sec:threadgroup-memory} and \cref{sec:imageblock-per-core-tile}.

\begin{xltabular}{\linewidth}{@{}>{\hsize=0.88\hsize}Xll>{\hsize=1.18\hsize}XX@{}}
\caption{Load and store forms used in this chapter.}\label{tab:memory-forms}\\
\toprule
Glossary name & Opcode & Lengths & Width and notes & Standing \\
\midrule
\endfirsthead
\tablecontinued{5}
\toprule
Glossary name & Opcode & Lengths & Width and notes & Standing \\
\midrule
\endhead
plain 32-bit load & \opn{12682} & 8, 10, 14 & 4 bytes, element code 4 & measured (MM~\mm{25.197}; below Metal) \\
16-bit load & \opn{12646} & 8, 10, 14 & 2 bytes (half, bf16, short, char); optional high-half destination & measured (MM~\mm{25.141.3}) \\
plain 64-bit load & \opn{12691} & 8, 14 & 8 bytes into a pair, element code 8 & measured (glossary) \\
three-component vector load & \opn{12700} & 8, 14 & 12 bytes into a 3-register tuple; the access steps 16 & measured (MM~\mm{25.90}) \\
128-bit vector load & \opn{12709} & 8, 10, 14 & 16 bytes into a 4-register tuple, element code 16; byte displacement & measured (MM~\mm{25.141.18}) \\
immediate-address one-word load & \opn{12688} & 8, 10, 14 & word at (index + displacement) $\times$ element; mask 2066 & measured (MM~\mm{25.177}) \\
tensor load, half/bf16 & \opn{12674} & 12, 16 & two words (four half/bf16), element code 2; immediate element mask & measured (MM~\mm{7}) \\
masked tensor load, half/bf16 & \opn{12675} & 12, 16 & as \opn{12674}, with a 16-bit mask register & measured (MM~\mm{7}) \\
int8 tensor load, one word & \opn{12656} & 12, 16 & one word (four int8), element code 1 & measured (MM~\mm{7}) \\
plain 32-bit store & \opn{17229} & 8, 10, 14 & 4 bytes (sub-form 00) & measured (MM~\mm{25.117}, MM~\mm{25.145.4}) \\
128-bit vector store & \opn{17256} & 8 (corpus) & 16 bytes from a 4-register tuple; uses every wait bit 24--31 & measured (MM~\mm{6}) \\
store, sub-form 01, texel/word-slot & \opn{17235} & 8, 10, 14 & one word at an immediate address; reads a texture fetch; mask 2066 & measured (MM~\mm{25.62}) \\
two-component store & \opn{17244} & 8, 10, 14 & words k and k+1 from a pair; mask 2098 & measured jointly with \op{17262} (MM
\mm{25.177}) \\
tensor store & \opn{17257} & 10, 12, 16 & four words; immediate mask & measured (MM~\mm{6}, MM~\mm{7}) \\
masked tensor store & \opn{17258} & 10, 12, 16 & four words; mask register, bit j enables word j & measured (MM~\mm{7}) \\
threadgroup-memory load & \opn{12364} & 8, 14 & 2 or 4 bytes & measured (\cref{sec:threadgroup-memory}) \\
threadgroup-memory store & \opn{13288} & 8, 10, 14 & 2 or 4 bytes & measured (\cref{sec:threadgroup-memory}) \\
imageblock load, 32-bit & \opn{12151} & 14 & one 32-bit member & measured (\cref{sec:imageblock-per-core-tile}) \\
imageblock store, 32-bit & \opn{13075} & 14 & one 32-bit member & measured (MM~\mm{25.96}) \\
\bottomrule
\end{xltabular}

Apple's compiler also emits \op{12673} to load pointer +
immediate from a register-pair base, \op{12073} for a runtime 32-bit address into thread-private arrays (MM~\mm{25.39.2}),
and the 64-bit store \op{17238} (no glossary entry; lengths 8 and 10), none of which is in \cref{tab:memory-forms}.
\op{12692} (element code 4, measured, MM~\mm{25.37}) and the masked twins of the int8 and fp32 tensor loads are also
omitted. The 16-bit store (\opn{17193}, lengths
10 and 14) is measured (glossary), and \op{17202} is measured in the packed int8 epilogue (MM~\mm{25.105}). The rest of the
immediate-address family is in \cref{sec:immediate-address-vector-forms}, the spill stores in \cref{ch:register-file}. Two
forms seen only in MLX's quantized matmul, \op{12352} and \op{13276}, have an open address space (MM~\mm{25.144.1}).

Decoder-only and unmodelled forms are in MM~\mm{25.88} and MM~\mm{25.192}. The repository's encoders round-trip 1,145 of
1,190 and 3,419 of 3,793 loads and stores; the unclaimed sub-forms are named in MM~\mm{15}.

In the decoder, 132 memory opcodes declare implicit uses of \code{SR\_LOCAL\_X} and \code{SR\_LOCAL\_Y} (\op{12118} to
\op{12135} among them, no glossary entry for either); none occurs in the shipped corpus
(\code{isa/g17-local-indexed-memory.txt}).

\section{Immediate-address and vector forms}\label{sec:immediate-address-vector-forms}

The immediate-address family's operands, in \op{17235}'s order, are value or destination, operand 1,
mask, binding, 0, index, displacement, element (MM~\mm{25.177}); the address is the word at \code{(index + displacement) * element} of the
binding, with index and displacement both immediates. The mask field is a component mask in bits 4--7 above 0x802:
2066 one word, 2098 two, 2162 three, 2290 four (bits 4--7 = 1, 3, 7, 15), filled from the tuple's lowest register up;
2065 is the halfword form (MM~\mm{25.195}). The family's members are \op{12688}, the immediate-address two-word
load (\opn{12697}, MM~\mm{25.184}), \op{12706}, \op{12715}, \op{17244}, \op{17253}, \op{17262} and \op{17199}.

The first model of this family (emulated, then executed on hardware; MM~\mm{25.177}) loaded each word with element size 1, keeping only its low byte, and still matched the GPU-output hash recorded
from Apple's build on 9 of 10 census sources, because the census inputs are small integers, (7i + 3) mod 61. Only the one source with
float data failed; it read as zero (5.0 is 0x40a00000). Rerun with full-width random buffers (full-mantissa floats with
exponents to $2^{\pm20}$, uints over all 32 bits), all 10 sources match Apple's GPU exactly, and each of three deformed
models is refuted for every form (\cref{tab:immediate-address-refutation}). The two-component store was confirmed only
jointly with the four-word store. The same protocol later confirmed the three-word store (6 of 6,
MM~\mm{25.192}) and the 16-bit store at an immediate address (3 of 3, MM~\mm{25.195}). A memory form is marked
measured in this reference only if its test data could have exposed a wrong width.

\begin{longtable}{@{}lllll@{}}
\caption{Immediate-address forms on full-width random data: sources that match Apple's GPU, and sources on which each
deformed model differs (MM~\mm{25.177}).}\label{tab:immediate-address-refutation}\\
\toprule
Form & Sources matching & Low-byte load & Reversed load & First-only store \\
\midrule
\endfirsthead
\tablecontinued{5}
\toprule
Form & Sources matching & Low-byte load & Reversed load & First-only store \\
\midrule
\endhead
one-word load & 5 & 5 & 3 & 3 \\
four-word load & 6 & 6 & 6 & 6 \\
four-word store & 8 & 6 & 6 & 8 \\
two-component store & 2 & 1 & 1 & 2 \\
\bottomrule
\end{longtable}

\op{12709} takes a byte displacement (\cref{sec:address-computation}). Apple's compile of MLX's qmv carries 16, 32, 48 for
\code{x[k+4]}, \code{x[k+8]}, \code{x[k+12]}. \code{asm.encode\_vec4} calls the field "units of eight bytes"; for loads it is
bytes (MM~\mm{25.141.18}). The store's displacement has been seen only at 0, 16 and 32, which does not fix its units.

Tensor loads and stores move popcount(mask) consecutive elements into the tuple, lowest address in the lowest half
(measured, \code{tools/g17emu.py} EVIDENCE{[}12674{]}). int8 operands need \op{12656}, not the
generic four-word load (measured, MM~\mm{7}).

This compiler folds the base of a vector placed past a 32,767-byte displacement into the index register once, and
two per-channel vectors must lie within one field of each other (measured, MM~\mm{25.165}). With loop-carried
addresses (\code{ptr\_addr}), later loads read one carried index at displacements 16, 32, 48 and instructions per trip fall
121 to 109, 184 to 171, 69 to 59 and 309 to 295 in four qmv loops, bit-exact against a 0x7f sentinel (MM~\mm{25.141.18}).

\section{Masked and predicated access}\label{sec:masks-predicated-memory}

\op{612} computes \code{bit j of dst = [lo <= src + j < lo + w]} for j in 0..3, unsigned. The second
immediate is a \emph{width}. It was first read as an upper bound because nearly every witness had lo = 0, where the two
readings coincide; two cases that separate the readings falsified the bound (measured, MM~\mm{7}). lo and w are 8-bit, so they cannot
span a whole buffer, and lowerings compute the predicate relative to the tile.

\op{621} gives a mask of \code{min(max(W - src, 0), 4)} low ones, W = operand 5, measured at
W = 16, 32, 48 and 64 (src W-3..W gives 7, 3, 1, 0); the cap of four did not move with operand 3 at 8, 16, 32, 56
(\code{isa/g17-address-formula.txt}). An earlier fit, \code{(0xFFFFFFFF >> src) \& 0xF}, describes only the W = 32 case.
Every dispatch behind it had been authored on one witness whose operand 5 was 32, which coincides with the fitted
saturation point.

In \op{12675} and \op{17258}, mask bit j enables word j of the lane's
four-word access, and bits above 3 are ignored. Unwritten words kept their contents, only the declared M $\times$ N region was
written, and poisoned padding was never written, in six patterns across 32 lanes (measured, MM~\mm{7}).

In ordinary predicated access (measured, \code{isa/g17-predication.json}), a masked-off load suppresses the \emph{value} (inactive
lanes keep their sentinel) and a masked-off store suppresses the \emph{access} (seven masks); effects are per lane and nested
masks intersect. Whether a masked-off load still issues its access is unknown: the intended observable was a fault,
and no out-of-range load faulted (\cref{sec:out-of-range-accesses-observed}). Apple's compiler emits predicated
memory branchlessly between exec-mask instructions, with a skip branch for two-armed effects (\cref{ch:control-flow-execution}).

\section{Out-of-range accesses}\label{sec:out-of-range-accesses-observed}

The masking experiment needed a load that faults, to separate a suppressed access from a suppressed value (measured,
\code{isa/g17-predication.json} \code{wild\_address}). Unconditional \op{12682} reads at
4,194,304, 67,108,864, 268,435,456 and 1,073,741,824 words past a 256-word buffer (16 MiB to 4 GiB) ran through Metal,
each in a fresh process so that a device loss would cost only that process. All completed with status 0 and returned 0,
and an in-bounds control returned its data.

Out-of-range store behaviour has not been measured. A timing driver once read and wrote up to 2.78~MB
outside its buffers, with uncharacterised effect (MM~\mm{17} rule 34). The emulator refuses a store outside a buffer; that
refusal is a limit of the model (\code{tools/g17emu.py}).

Through Metal, one BIF0 read page fault came from a \code{regprobe\_multi} tensor-kernel probe whose configuration was not
recorded; a different process, sized from a sibling kernel's layout, hung the GPU for about 320 s and the machine
rebooted (MM~\mm{17} rule 34, as corrected; \cref{app:a}).

Below Metal, malformed launch state produced BIF0 read page faults at unmapped addresses while Submit still returned
0, both callbacks ran and the recovery counter advanced: a guessed 2,048-byte threadgroup resource word (fault at
address 0; the matched word then ran), an all-zero allocation 0 (fault at 0x10000010000, its end), an oversized
allocation layout (0x100000d0000, 0x100000d8000) and a malformed code-placement packet (0x80) (measured,
\code{docs/g17-first-dispatch-trials.md}). The requester of those reads was not identified. MM~\mm{17} rule 34's "only the
accelerator's own reads have faulted" therefore applies only to work submitted through Metal.

Undersizing a host dispatch's completion scope (\code{dim\textasciicircum{}2 * esz})
returns wrong or zero data past that offset with no error, which looks like a range error (measured, MM~\mm{17} rule 5). A dead load landing late can
overwrite a reused address register so that the next load reads out of range and returns 0 (\cref{sec:dependences-issue-scalar}, MM~\mm{25.114}).

\section{Threadgroup memory}\label{sec:threadgroup-memory}

Apple's compiler refuses more than 32,768 bytes per threadgroup: 48, 60 and 64 KiB are refused and 96 KiB drops the
compiler connection (measured, MM~\mm{7}). Metal itself does not check the limit for an authored container: a container
declaring 1 MiB of static threadgroup memory built a pipeline that reports \code{staticThreadgroupMemoryLength = 1048576}
(measured, MM~\mm{25.147}, nothing dispatched). What the hardware does with an oversize declaration is unmeasured, and
the emulator refuses threadgroup accesses past 32,768 bytes.

Per core the budget is 57,344 bytes (56 KiB), inferred from a step in run time when two threadgroups no longer fit
(measured step, MM~\mm{25.49}; the step and the residency rules are in \cref{ch:device-execution-model}).
\code{maxTotalThreadsPerThreadgroup} stays 1,024 across threadgroup memory from 64 B to 32,768 B (measured, MM~\mm{25.47}),
and \code{staticThreadgroupMemoryLength} is the value the kernel's metadata declares (MM~\mm{25.147}).

Threadgroup memory is written by \op{13288} and read by \op{12364}; Apple's compiler also emits a constant-index threadgroup
load, \op{12367} (no glossary entry; \code{isa/g17-threadgroup-pair.txt}). Their operands are value or
destination, operand 1, width, base, 0, index, index lifetime, displacement, element. The base's binding expression,
\code{bin(op0, const(K), 2)}, has scale 2 where device forms use 8, with 0 exceptions in 25,553 corpus memory instructions
(decoder, \code{isa/g17-address-formula.txt}). The address is base + index $\times$ element + displacement (the emulator's model,
checked against its recorded hardware runs, \code{tools/g17emu.py} EVIDENCE{[}13288{]}).

The listing shows a store and the load of the same address, compiled together in one kernel rather than paired up
afterwards (compiled by Apple and decoded,
\code{isa/g17-threadgroup-pair.txt}, \code{a6-tgidx}; the same shape executed below Metal).

\begin{lstlisting}
threadgroup-memory store (op13288), 10 bytes  0f040302030611c00020
    operands: R105 (the stored value), operand 1 = 0x80000010, 2066, bin(op0,const(4),2), 0,
              index register, 16, 0, 4      (op0: uniform register file)
threadgroup-memory load (op12364), 14 bytes   0f0403021a2610c0410080000000
    operands: R105 (the destination), operand 1 = 0x2001800000, 2066, bin(op0,const(4),2), 0,
              index register, 16, 0, 4
\end{lstlisting}

The store's operand 1 is 0x80000010: bit 31 set, a wait on slot 7, the slot this kernel's loads fill (\cref{sec:scoreboard}); the
trailing 4 is the element width. Below Metal, a program of this shape, each lane storing \code{A[i]} to scratch, a
threadgroup barrier (\cref{sec:barriers-fences}), then lane i reading slot \code{i XOR 32}, produced 64 of 64 values \code{11 * (i XOR 32) + 5} across two
simdgroups, and a 2,048-byte allocation held 512 separately addressable words with every written value recovered
(measured, \code{docs/g17-first-dispatch-trials.md}, "Parallel static-threadgroup launch mapping"). The static size is a
launch resource word, not an instruction field (\cref{ch:submission-below-metal}): a guessed value for 2,048 bytes faulted, and the matched
value Metal writes (\code{0f00080c00000000}) does not equal size/8, as it does for the smaller sizes.

An ordinary threadgroup load or store provisions the threadgroup region: removing the one threadgroup store
made a threadgroup atomic write nothing (measured, ledger \code{g17-the-threadgroup-atomic-subtracts-and-needs-a-region}).

In the threadgroup load, clearing byte4{[}3{]}, byte4{[}4{]} or byte6{[}4{]} made the read return 0;
byte5{[}5{]}, byte8{[}6{]} and byte10{[}7{]} are inert; of the three-bit variants, 000 read correctly and 010, 100 and 110 read 0
(measured, \code{isa/g17-execution-tgbits-results.json}, \code{-tgtrio-results.json}). These are the positions of the 128-bit
vector load's fill-slot field (MM~\mm{6}); whether they are a slot field here is not established.

\section{Imageblock (per-core tile memory)}\label{sec:imageblock-per-core-tile}

The imageblock is a per-threadgroup tile indexed by thread position and shared within the threadgroup (measured,
\code{isa/g17-imageblock-execution.txt}): 32 lanes wrote 32 distinct values to their own pixels, and every lane that read pixel
(0, 0) wrote lane 0's value into a buffer prefilled with 0xdeadbeef.

The imageblock store \op{13075} and load \op{12151} take the operands
\code{[value, operand 1, width 18, member byte offset, 1, coordinate register, coordinate lifetime, dx, dy]} (measured,
MM~\mm{25.145.2}). The coordinate register packs x in its low 16 bits and y in its high 16; Apple's compiler builds it with
\op{10286}. The declared pipeline allocated 128 bytes for 32 lanes of one 32-bit member
(MM~\mm{25.104.1}). The emulator refuses coordinates outside the tile, although one imageblock read at
x = 32, outside a 32-wide tile, completed on hardware (it was not scored; MM~\mm{25.104.2}).

Finding the imageblock store's wait bit took eleven rounds on hardware (MM~\mm{25.96}, MM~\mm{25.104}). The target was a
tensor GEMM, every byte from this compiler, that stages its output through the imageblock. In every round a neighbour
read returned the reader's own value or 0, and three separate defects produced that result.

First, the store released its coordinate register (operand 6 = 16) while the following read
reused it (release semantics in \cref{ch:register-file}). A one-bit same-run control, keep (0) against release (16), collapsed
all 32 lanes to element 0 (measured, MM~\mm{6}; \code{isa/g17-imageblock-dseries-receipt.json}), and the fix was to set
operand 6 from liveness.

With the coordinate kept, the x offset (operand 7 = 1) still had no visible effect in a tensor
program. Apple's compiler never uses it for a neighbour. It computes \code{(tid.x + 1) \& 31} with \op{10289} and
\op{435} and reads with operand 7 = 0, and this compiler addresses the neighbour the same way, by an explicit column
(measured, MM~\mm{25.104.2}).

Apple's tensor kernel then read its neighbour on all 32 lanes and this compiler's did not. The only
field-level difference was the store's operand 1, Apple's 0x80000010 against this compiler's 0x01000010. Swapping it
alone moved the result both ways, and bit 31 decided: 0x80000010 and 0x81000010 passed, 0x00000010 and 0x01000010
failed. The reading recorded at the time, that bit 31 selects storage other lanes can read, was later refuted
(\cref{app:a}; MM~\mm{25.104.3}). Bit 31 is the store's wait on slot 7.

Bits 24--31 of the imageblock store's operand 1 are its wait mask, the same field general stores and MMAs carry
(\cref{sec:scoreboard}). In a census of all 57 corpus
imageblock-store instances the bits equal exactly the slots the store's sources are pending on: 47 need none and carry none, 8 carry bit
24 for the coordinate's special-register reads, 2 carry bit 25. A prediction made before compiling, that an Apple kernel
storing a freshly loaded value (its load fills slot 7) would carry bit 31, was confirmed. \Cref{tab:imageblock-wait-arms} gives
the hardware runs of that kernel (\code{h[t] = 5t + 3}, a device barrier, \code{v = h[t \textasciicircum{} 1]}, an own-pixel store of v, read back), three per arm.

\begin{xltabular}{\linewidth}{@{}>{\hsize=1.19\hsize}X>{\hsize=0.81\hsize}X@{}}
\caption{Imageblock store operand 1 against lanes that read their neighbour's value correctly, in Apple's kernel that stores a
freshly loaded value. Three runs per arm.}\label{tab:imageblock-wait-arms}\\
\toprule
store operand 1 & lanes correct \\
\midrule
\endfirsthead
\tablecontinued{2}
\toprule
store operand 1 & lanes correct \\
\midrule
\endhead
0x80000010, as compiled (waits slot 7) & 32 / 32 \\*
0x10, wait cleared & 0 / 32, every lane 0 \\*
0x2000010, waits slot 1, which nothing fills & 0 / 32, every lane 0 \\*
\bottomrule
\end{xltabular}

Without the right slot the store takes its value register before the load writes it. Earlier deformations of operand 1
in Apple kernels had found the bits inert because those kernels stored
ALU results, which need no wait. Round 11 routed the stored value through a waiting copy and kept the slot-0 wait for
the coordinate (0x01000010); every arm was bit-exact (MM~\mm{25.104.4}). Bit 31 does not decide whether other lanes
can read a word, since the imageblock is shared within the threadgroup whatever bit 31 holds (above). The emulator
still allows cross-lane reads only of words stored with bit 31; no measurement requires that restriction.

An undeclared imageblock reads 0 even when the tile is sized at encode time with \code{setImageblockWidth} (measured,
MM~\mm{25.104}). The test kernel stored each read plus 1.0, and C{[}0, 0..31{]} held 1.0 in each of three runs, so every read
returned 0. The pipeline must come from \code{newComputePipelineStateWithDescriptor:}. A pipeline looked up by function
name does not carry what the driver needs to size the tile, and \code{imageblockMemoryLengthForDimensions:} then reports
768 bytes at 32 $\times$ 1. \code{setImageblockWidth:height:} must match the threadgroup size
(\code{isa/g17-imageblock-execution.txt}; declaring metadata in \cref{ch:kernel-metadata-image}).

The imageblock barrier is \op{447} at scope 532, bytes \code{475100060000}
(\code{threadgroup\_barrier(mem\_threadgroup\_imageblock)}), which Apple's own-element kernel carries; Apple's compiler elided a
source \code{simdgroup\_barrier(mem\_threadgroup)} between a one-simdgroup imageblock store and load (Apple's compiler emits,
\code{isa/g17-barrier-scopes.txt}, \code{isa/g17-imageblock-execution.txt}).

Use by more than one simdgroup or threadgroup has never been witnessed (MM~\mm{25.145.2}). Staging a
fragment through the imageblock is bit-exact but frees no registers; its cost is in \cref{sec:chaining-feed-modes} (MM~\mm{25.106}).

\section{Constants and uniforms}\label{sec:constants-uniforms}

A constant is an immediate inside the instruction, an entry in the \emph{constant pool}
(compile-time data in metadata slot 13 of the container; Apple's compiler emits, \code{agxforge/g17/cooperativemetadata.py}), or a
\emph{published uniform}, a value Apple's compiler computes per dispatch in a separate \emph{constant program} that runs once before the main
program and writes the uniform register file (binary format of the program and the pool in \cref{ch:kernel-metadata-image}). Main reads pool entries and published uniforms through the same operand
syntax, \code{bin(op0, const(K), W)}, so the operand alone does not say which one it is. Treating a published uniform as a
pool entry gave 6 silent wrong answers (measured, MM~\mm{25.179}).

The emulator reproduces Apple's GPU output by running the constant program first, as one thread over the same buffers,
and handing its uniform writes to main: 6 of 6 sources are byte-identical on full-width random buffers, and three deformed
models (every published uniform reads 0; the pointer names the next binding; a 32-bit uniform lands with halfwords
swapped) each differ on all 6 (measured, MM~\mm{25.179}). In constant programs the base is a register pair holding a
binding pointer that \op{14061} fetched, and the address is pointer plus byte offset (measured, MM~\mm{25.179}). The
binding pointers the constant program reads and its
uniform-destination writer forms are given in \cref{sec:constant-program-record}; a uniform destination must be written
with one value across the active lanes. Of 15,256 corpus constant-program uniform writes, 14,814 are read by main; the
442 never-read writes are all uniform / staging writes, and why main never reads them is not known (decoder, MM~\mm{25.120}).

In main, \op{428} takes its mask from a pool operand (\cref{sec:integer-arithmetic}). The replay of the shipped
decode, which uses \op{428} 768 times, refused until it loaded the pool (measured, MM~\mm{25.197}). \op{435} takes an 8-bit immediate; a wider mask such as 0xFF0 goes to the
uniform-operand \op{437} (Apple's compiler emits, MM~\mm{25.192}).

In the uniform / staging write (decoder, MM~\mm{25.120.1}), the value the decoder printed as a
"publish index" is a scoreboard slot + 1. With a uniform destination it is 0 (47,269 of 47,269); with a destination in the message staging file
(\opn{4}) it is 1--8 (131,313 of 131,313). The reader of a staging write is a message instruction (\op{14661}, \op{10909}, \op{15813}), and it waits on slot index - 1: of 189,898 staging writes,
174,909 are waited on by the reader and 14,989 by an intervening instruction, so every one is waited on. A staging write is a
late value, like a load (\cref{sec:scoreboard}).

\section{Textures}\label{sec:textures}

Textures are read through \op{15813}, \op{14661}, \op{10909} and \op{11564} (glossary), and no operand of these carries the coordinate. Two reads at different coordinates are byte-identical, and the
coordinate is published into the message staging file (\opn{4}) at \code{[op4+0*4]} and \code{[op4+2*4]} before the message (measured,
\code{tools/g17texrun.py}).

An authored read, \code{A[400] = TX.read(x, y).x}, was executed (\code{tools/g17texrun.py},
\code{isa/g17-execution-texture-results.json}) on a texture holding y*1000 + x + 7, chosen so a wrong texel, a wrong row, a broadcast coordinate and an unbound texture
give four different wrong answers: coordinate (5, 1) gave 1012, (2, 3) 3009, and lane 31 alone at (7, 3) 3014; all 32
lanes racing for one word gave an ambiguous 3007. The same runs showed that the texture operand is a dense index over
the textures the function uses, not the Metal binding index (five constructed kernels); and publishing a register
coordinate with a uniform / staging write needs operand 2 = 1048576; with 16777216 every lane silently reads texel (0, 0).

\op{17235} reads a texture fetch (measured,
MM~\mm{25.62}), and the 14-byte two-component store \opn{17244} read one once (texel 201 at (1, 2), one 4 $\times$ 4 R32Uint texture, one
thread; MM~\mm{25.146}). Whether the plain 32-bit store \op{17229} can is open.

A dispatched bisection of Apple's reading program found that swapping only its store for this compiler's 8-byte plain
store stopped the read. That store, and a 12-byte ALU copy placed after the fetch, saw the destination register's
previous contents, at zero and at six intervening instructions, and the ledger concluded "not timing" (ledger
\code{g17-only-one-consumer-can-read-a-texture-fetch}). Apple's \code{tex2d-read}, however, stores the texture read's
result with an 8-byte plain store whose operand 1 is 16777232 (0x1000010, a wait on slot 0), and that read publishes on
slot 0 (Apple's compiler emits, MM~\mm{25.146}, \code{tools/g17applestorefetch.py}). This compiler's encoder writes no
wait on the 8-byte plain store, and the ALU copy in the bisection waited on slot 7. The bisection's arms are therefore
equally consistent with a missing wait on the fetch's slot. "Not timing" does not rule that out, since unwaited late
values stayed wrong after 64 intervening instructions elsewhere (\cref{sec:scoreboard}). The settling dispatch, Apple's
\code{tex2d-read} run and read back, has not been made.

Pixel format is descriptor-side (Apple's compiler emits, \code{isa/g17-texture-sampler-surface.json}):
\code{texture2d<uint>} and \code{texture2d<float>} (R32Uint, R32Float) compile to identical code and metadata; R16Float changes
only the instructions that operate on half. A sampler adds one register and a 4-byte sampler-state descriptor and swaps the
texture read for the texture sample. Texture metadata slots are listed in \cref{ch:kernel-metadata-image}. No sampled or write texture has been dispatched.

\textbf{Not known.} The vector store's displacement units; the source of \op{621}'s cap of four, which is not a four-bit field of the
encoding; out-of-range loads
of other forms, resource kinds and offsets; device accesses below Metal at arbitrary addresses or in other allocation
classes; the bit-4 flag (0x10) every imageblock store's operand 1 carries (MM~\mm{25.96}), and the imageblock read's
operand-1 word Apple emits, 0x2000800000 (MM~\mm{25.104.2}).

\textbf{Evidence.} MM~\mm{0.8}, MM~\mm{6}, MM~\mm{7}, MM~\mm{15}, MM~\mm{16}, MM~\mm{17} rules 5, 6 and 34, MM~\mm{25.47}, MM~\mm{25.49}, MM~\mm{25.62}, MM~\mm{25.89}, MM~\mm{25.90},
MM~\mm{25.96}, MM~\mm{25.104} (25.104.1 to 25.104.4), MM~\mm{25.105}, MM~\mm{25.106}, MM~\mm{25.114}, MM~\mm{25.114.3}, MM~\mm{25.120}, MM~\mm{25.120.1},
MM~\mm{25.141.18}, MM~\mm{25.144.1}, MM~\mm{25.145.2}, MM~\mm{25.146}, MM~\mm{25.147}, MM~\mm{25.165}, MM~\mm{25.177}, MM~\mm{25.179}, MM~\mm{25.184}, MM~\mm{25.192},
MM~\mm{25.195}, MM~\mm{25.197}; \code{isa/g17-address-formula.txt}, \code{isa/g17-addressing-rules.json}, \code{isa/g17-predication.json},
\code{isa/g17-threadgroup-pair.txt}, \code{isa/g17-imageblock-execution.txt}, \code{isa/g17-barrier-scopes.txt},
\code{isa/g17-wait-mask-census.json}, \code{isa/g17-texture-sampler-surface.json}, \code{docs/g17-first-dispatch-trials.md},
\code{tools/g17emu.py}, \code{tools/g17texrun.py}, \code{tools/g17applestorefetch.py}.
